%=========== ΛΥΣΗ - 14ο ΘΕΜΑ Δ ===========
\begin{thema}{Δ}
\begin{erwthma}
\item Για κάθε $x>0$ έχουμε
\begin{align*}
f'(x)
&=\frac{\ln x}{x^2}\\
&=-\ln x\left(\frac{1}{x}\right)'\\
&=-\left(\frac{\ln x}{x}\right)'
+\frac{1}{x}(\ln x)'\\
&=-\left(\frac{\ln x}{x}\right)'
+\frac{1}{x^2}\\
&=\left(-\frac{\ln x}{x}-\frac{1}{x}\right)'\\
&=\left(-\frac{\ln x+1}{x}\right)'.
\end{align*}
Σύμφωνα με τις συνέπειες του Θεωρήματος Μέσης Τιμής, υπάρχει σταθερά $c\in\mathbb{R}$ τέτοια ώστε
\[
f(x)=-\frac{\ln x+1}{x}+c.
\]
Από την αρχική συνθήκη $f(1)=0$ παίρνουμε
\[
-1+c=0\Leftrightarrow c=1.
\]
Επομένως
\[
{f(x)=1-\frac{\ln x+1}{x}},\quad x>0.
\]

\item Επειδή $x^2>0$, το πρόσημο της
\[
f'(x)=\frac{\ln x}{x^2}
\]
είναι το πρόσημο του $\ln x$. Άρα:
\begin{itemize}
\item $f'(x)<0$ για $0<x<1$,
\item $f'(1)=0$,
\item $f'(x)>0$ για $x>1$.
\end{itemize}
Στον παρακάτω πίνακα βλέπουμε το πρόσημο της $f'$ και τη μονοτονία της $f$:
\begin{center}
\begin{tikzpicture}
\tkzTabInit[color,lgt=2,espcl=1.8,colorC = maincolor!50!white,colorL = maincolor!20!white,colorV = maincolor!50!white]%
{$x$ / .8,$f'(x)$ /0.8,$f(x)$/0.8}%
{$0$,$1$,$+\infty$}%
\tkzTabLine{d, -, z, +, }
\tkzTabVar{D+/{+\infty},-/{0},+/{1}}
\end{tikzpicture}
\end{center}

Άρα η $f$ είναι \textbf{γνησίως φθίνουσα} στο $(0,1]$ και \textbf{γνησίως αύξουσα} στο $[1,+\infty)$. Επομένως παρουσιάζει ολικό ελάχιστο στο $x_0=1$, με ελάχιστη τιμή
\[
{f(1)=0}.
\]

\item Η $f$ είναι συνεχής και γνησίως φθίνουσα στο $(0,1]$. Επιπλέον,
\[
\lim_{x\to0^+}f(x)=+\infty
\quad\text{και}\quad
f(1)=0<\frac{1}{2}.
\]
Επομένως, από το Θεώρημα Ενδιάμεσων Τιμών και τη γνήσια μονοτονία της, η εξίσωση
\[
f(x)=\frac{1}{2}
\]
έχει ακριβώς μία ρίζα στο $(0,1)$.

Στο $[1,+\infty)$ η $f$ είναι γνησίως αύξουσα. Επίσης
\[
f(e)=1-\frac{2}{e}<\frac{1}{2},
\]
αφού $e<4$, ενώ
\[
f(e^2)=1-\frac{3}{e^2}>\frac{1}{2},
\]
αφού $e^2>6$. Άρα, από το Θεώρημα Ενδιάμεσων Τιμών, υπάρχει μία ρίζα στο $(e,e^2)$, η οποία είναι μοναδική λόγω της γνήσιας αύξησης της $f$.

Συνεπώς η εξίσωση $f(x)=\frac{1}{2}$ έχει ακριβώς δύο ρίζες στο $(0,+\infty)$.

\item Από τον τύπο της $f$ έχουμε
\begin{align*}
\int_1^e f(x)\d x
&=\int_1^e\left(1-\frac{\ln x+1}{x}\right)\d x\\
&=\left[x-\frac{(\ln x)^2}{2}-\ln x\right]_1^e\\
&=\left(e-\frac{1}{2}-1\right)-1\\
&={e-\frac{5}{2}}.
\end{align*}
\end{erwthma}
\end{thema}
%=========== ΤΕΛΟΣ ΛΥΣΗΣ - 14ο ΘΕΜΑ Δ ===========
